\subsection{An explicit counterexample on the cyclotomic tripod}
Put $F=\Q(\mu_\infty)$. We use the valuation calculation underlying R18,
and pass through an explicit degree-nine cover of the tripod. This
answers the tripod assertion in C1, with all maps specified.

\begin{statement}{C.1}{Proved}{Projectivity over the full cyclotomic field}
The profinite group $G_F$ is projective in the category of profinite
groups and continuous homomorphisms. Every continuous surjection of a
profinite group onto $G_F$ has a continuous section.
\end{statement}
\begin{proof}
Let $L/F$ be finite. A central simple algebra over $L$ descends to a
number field $k\subset L$, because it is defined by finitely many
structure constants. Milne \cite[VII, Proposition 7.2]{MilneCFT}
states that every Brauer class of a number field is killed by a finite
cyclotomic extension of that number field. Since all roots of unity
belong to $L$, we obtain $\operatorname{Br}(L)=0$.
This verifies the projectivity input used in R18.

Fix a prime $\ell$ and a pro-$\ell$ Sylow subgroup $P\subset G_F$.
Its fixed field $F_P$ is algebraic over $F$. Every central simple
algebra over $F_P$ descends to a finite intermediate field, so
$\operatorname{Br}(F_P)=0$. Since $\mu_\ell\subset F$, Kummer theory gives
\[
 H^2(P,\mathbb F_\ell)
 =H^2(G_{F_P},\mu_\ell)
 =\operatorname{Br}(F_P)[\ell]=0.
\]
A finite $\ell$-primary continuous $P$-module has a composition series
with trivial $\mathbb F_\ell$ factors: the action factors through a
finite $\ell$-group, whose only simple modules in characteristic
$\ell$ are trivial. The long exact sequence of cohomology gives
$H^2(P,M)=0$ for each such module.

Restriction $H^2(G_F,M)\to H^2(P,M)$ is injective for every finite
$\ell$-primary $G_F$-module $M$. A class that becomes zero on $P$
becomes zero on an open subgroup $H\supset P$: a cobounding continuous
cochain has finite image, and its identities hold on a sufficiently
small open neighbourhood of $P$. Corestriction composed with
restriction multiplies the class by $[G_F:H]$, prime to $\ell$, so
the class is zero. The transfer formula is
\cite[II, Proposition 1.30]{MilneCFT}; it applies here through the
finite quotients carrying the cochains.

Thus $H^2(G_F,M)=0$ for every finite discrete torsion module.
Every discrete torsion module is a filtered union of finite invariant
submodules: each orbit is finite, and finitely many finite-order
elements generate a finite abelian group. Continuous cochains commute
with these unions. Dimension shifting gives vanishing in all degrees
greater than one, hence $\operatorname{cd}(G_F)\leq1$.
Serre \cite[I, \S5.9, Corollary 2 to Proposition 45, p.~58]{SerreGC}
identifies cohomological dimension at most one with projectivity for
an arbitrary profinite group. This gives the stated lifting property.
No local section theorem is used.
\end{proof}

\begin{statement}{C.2}{Proved}{A degree-nine cover with no rational compactification point}
The smooth plane cubic
\[
 D:\quad X^3+2Y^3+4Z^3=0
\]
has $D(F)=\varnothing$. The morphism $t=-4Z^3/X^3$ to $\Pj^1_F$
restricts to a geometrically connected finite \'etale degree-nine
cover $V\to U_F$. Its affine equations are
\[
 t=-4a^3,\qquad 1-t=-2b^3,\qquad a=Z/X,\quad b=Y/X.
\]
\end{statement}
\begin{proof}
A proposed $F$-point on $D$ is defined over some $\Q(\mu_N)$.
At a place over 2 the ramification index $e$ is a power of 2:
roots of unity of odd order give unramified extensions of $\Q_2$,
and $2$-power roots of unity give ramification degrees that are
powers of 2. Normalize the valuation to have value group $\Z$.
The three nonzero terms in the cubic have valuations congruent to
$0,e,2e$ modulo 3, respectively. These residues are distinct.
The minimum valuation therefore occurs exactly once, preventing
the sum from being zero. Zero coordinates do not affect this
argument. This verifies the explicit instance $p=2$, $a=b=c=1$
of Clark \cite[Theorem 1]{Clark2006}.

The partial derivatives of the cubic have no simultaneous projective
zero, so $D$ is smooth. A reducible plane cubic over an algebraically
closed field has intersecting components and hence a singular point.
Thus $D$ is geometrically integral. On $XYZ\ne0$, its coordinate
algebra over $F[t,t^{-1},(1-t)^{-1}]$ is obtained by the two displayed
cube-root equations. After multiplying by nonzero scalars they
are monic, and their derivatives in $a,b$ are invertible on this
open. This gives a finite \'etale algebra of rank $3\cdot3=9$.
It is the algebra of the geometrically integral open $V\subset D$.
\end{proof}

\begin{statement}{C.3}{Refuted}{Conjecture A for the tripod over $F$}
There exists a continuous section of $\pi_1(U_F)\to G_F$ that is
neither the section of an $F$-point nor cuspidal. Its mod-3
Kummer coordinates are
\[
 (A,B)=(2k,k),\qquad
 k=\kappa_3(2)\ne0\quad\hbox{in }H^1(G_F,\mu_3).
\]
\end{statement}
\begin{proof}
By C.1 choose a continuous section
$\sigma:G_F\to\pi_1(V)$. The finite \'etale morphism in C.2 induces
a continuous open injection
$\iota:\pi_1(V)\hookrightarrow\pi_1(U_F)$, after choosing compatible
geometric base points. The composite $s=\iota\sigma$ is a section.
Changing the chosen paths conjugates it by the geometric fundamental
group, precisely the ambiguity allowed here.

On $V$, the functions $t,1-t$ are $-4,-2$ times cubes.
Pullback along $\sigma$ of their Kummer classes gives
$A=\kappa_3(-4)$ and $B=\kappa_3(-2)$. Roots of unity are divisible
in $F^\times$, so this pair is $(2k,k)$. A cube root of 2 in a finite
cyclotomic field would give $3v(z)=e$ at a place over 2, with
$3\nmid e$. Thus $k\ne0$, and it has order 3.

For a point section $s_x$, the same coordinates are
$(\kappa_3(x),\kappa_3(1-x))$. Equality with $(2k,k)$ would give
$x=-4a^3$ and $1-x=-2b^3$ for $a,b\in F^\times$. Then
$(1:b:a)\in D(F)$, contradicting C.2. The comparison uses the
Kummer isomorphism
$H^1(G_F,\mu_3)=F^\times/(F^\times)^3$.

At cusp 0 the regular unit $1-t$ has value 1; every section in
that packet has $B=0$. At cusp 1 the analogous condition is $A=0$.
At infinity, $(1-t)/t$ is a regular unit of value $-1$, whose
Kummer class over $F$ is zero, so $A=B$. These identities follow
by restriction to the corresponding punctured henselian local
rings and decomposition groups. The pair $(2k,k)$ satisfies none
of these three cusp conditions. Hence $s$ is noncuspidal.

The finite mod-3 certificate is checked in
\texttt{code/root\_checks.py}. The existence of the continuous
splitting is proved by C.1; the program does not purport to
compute that noncanonical splitting.
\end{proof}

\begin{statement}{C.4}{Proved}{Birational lifting over the same field}
Every continuous section of $\pi_1(U_F)\to G_F$, including the
section of C.3, lifts to a continuous section of
$G_{F(t)}\to G_F$.
\end{statement}
\begin{proof}
For a representative $s$ form the profinite fibre product
\[
 P_s=\{(\gamma,g)\in G_{F(t)}\times G_F:q(\gamma)=s(g)\},
\]
where $q:G_{F(t)}\twoheadrightarrow\pi_1(U_F)$ is the continuous
quotient obtained by retaining extensions unramified over $U_F$.
The projection $P_s\to G_F$ is a continuous surjection.
By C.1 it has a continuous splitting. Composing with
$P_s\to G_{F(t)}$ gives the required lift, whose composite with
$q$ is the prescribed $s$.
\end{proof}
The field in C.3--C.4 is exactly $F=\Q(\mu_\infty)$.
These results supply no counterexample over $\Q$.
